\scanpage{34}
% archive sheet number: 33
% @ text / title, two lines
Détermination du module de Dieudonné filtré du groupe de Lubin-Tate.
% @ text / field and valuation ring
$K$ extension finie de $\mathbb Q_p$, anneau d'entiers $A$, corps résiduel
% @ diagram / square and data on its right
\diagram{\begin{tikzcd}[ampersand replacement=\&]
\mathbb Z_p\arrow[r,dash]\arrow[d,dash]\&A\arrow[d,dash]\& k\simeq\mathbb F_q,\ q=p^f\\
\mathbb Q_p\arrow[r,dash]\&K\&{[K:\mathbb Q_p]=h=ef}\\
\&\&\pi\in A\quad\text{uniformisante}
\end{tikzcd}}
% @ text / Witt vectors
$W=W_\infty(k)$ vecteurs de Witt,
% @ diagram / two inclusions with parenthesized degrees
\diagram{\begin{tikzcd}[ampersand replacement=\&]
\mathbb Z_p\arrow[r,hook,"{(f)}"']\& W\arrow[r,hook,"i","{(e)}"']\& A
\end{tikzcd}}
% @ text / aim, with two words entered above the line
On veut déterminer \inserted{d'abord} les groupes de B.T. \inserted{$G_0$} de hauteur $(=\unsure{\mathrm{rang}|v.\,\mathrm{g}})$ $h$ sur $\deleted{\unsure{\mathbb F_q|A|k}}\,k$, avec opération de $A$ dessus [telle que $\pi_{G_0}=F^f$ (automorphisme de Frobenius de $G_0/\mathbb F_q$)]
% @ text / Dieudonné theory, module and endomorphism
Par la théorie de Dieudonné, ils correspondent aux modules libres \inserted{$M$} de rang $h$ sur $W$, avec un endom $\sigma$-linéaire ($\sigma$ autom de Frob. de $W$) $F_M$, et opération de $A$ sur $M$ $(\lambda\mapsto\lambda_M)$ par $W$-\unsure{autom|endom} commutant à $F_M$, satisfaisant
% @ diagram / square on the left of the module conditions
\diagram{\begin{tikzcd}[ampersand replacement=\&]
A\arrow[r,dash]\arrow[d,dash]\& A\otimes_{\mathbb Z_p}W\arrow[d,dash]\\
\mathbb Z_p\arrow[r,dash]\&W
\end{tikzcd}}
% @ display / two starred requirements
\[
\left[(*)\qquad\deleted{\unsure{\pi_M|\lambda_M|\lambda_\pi}}\ F_M^f=\pi_{\underline M}\right]
\qquad
\begin{aligned}
(*')\quad&\exists\,V_M\ \text{t.q.}\\
&F_MV_M=V_MF_M=p
\end{aligned}
\]
% @ text / the two module structures
Les deux structures de module sur $M$ correspondent à une structure de module sur $A_W=A\otimes_{\mathbb Z_p}W$. Soit $\deleted{\unsure{F_{\mathbb Z_p}|F_W|\sigma_W}}\,\sigma_A$ l'autom de $A_W=A\otimes_{\mathbb Z_p}W$ donné par
% @ display / action on a simple tensor
\[
\sigma_A(\lambda\otimes w)=\lambda\otimes\sigma w
\]
% @ text / last two lines
Alors la donnée de $F_M$ équivaut à la donnée d'un endom $\sigma_A$-linéaire, satisfaisant $(*)\ (*')$.
